\subsection{Berechnung der Off-Label-Proportion}
\label{oluprop_method}

Wir interessieren uns für den Anteil der Anwendungsdauer der Medikamentenverordnungen, die Off-Label sind, im Verhältnis zur gesamten Anwendungsdauer der Medikamentenverordnungen. 
Diesen Anteil wollen wir für die verschiedenen Einrichtungen $k$ mit $k = 1,\dots, K$ und Diagnosen $d$ mit $d = 1,\dots, D$ vergleichen. Wir müssen diesen Anteil als
Merkmal eines Patienten $p$ mit $p = 1,\dots, P$ betrachten und nicht als Merkmal
einer Verordnung, da für einzelne Verordnungen dieser Anteil in der Regel $0$ oder
$1$ ist. Außerdem können wir nicht über eine Einrichtung oder Diagnose
aggregieren, da die Verordnungen in den Einrichtungen und Diagnosen nicht als
unabhängig voneinander angenommen werden können. Das kommt daher, dass ein Patient
meistens mehrere Verordnungen erhält und es anzunehmen ist, dass diese
Verordnungen miteinander zusammenhängen. Insbesondere ändert sich der Anteil der
Off-Label-Tage an der Anwendungsdauer mit der Anzahl der Verordnungen, die ein
Patient erhält, siehe Anhang A.1 Abbildung \ref{fig:OLU_Proportion_Patient_n}.\bigskip

Daher müssen wir für jeden Patienten $p$ den Off-Label-Anteil der Anwendungsdauer berechnen und können dann die Mittelwerte dieser Anteile
für die verschiedenen Einrichtungen und Diagnosen vergleichen. Wir bezeichnen diesen
Anteil im Folgenden als OLU-Proportion. Die OLU-Proportion eines Patienten $p$ ist
definiert als
\begin{equation}
    \label{eq:oluprop_patient}
    \text{OLU\_prop}_{p} := \frac{\sum_{i=1}^{n_p} \text{OLU\_Tage}_{p, i}}{\sum_{i=1}^{n_p} \text{Anwendungsdauer}_{p, i}}, p = 1,\dots, P,
\end{equation}
wobei $n_p$ die Anzahl der Verordnungen des Patienten $p$ ist. Die
OLU-Proportion ist somit ein Merkmal auf Patientenebene, wodurch wir uns auch
auf der typischen Einheit im medizinischen Kontext befinden.\bigskip

Wir könnten uns alternativ auch die durchschnittlichen Off-Label-Tage pro Patient
anschauen und zwischen den Einrichtungen und Diagnosen vergleichen. Dies würde
jedoch nicht berücksichtigen, dass die Off-Label-Tage direkt von der Anwendungsdauer
abhängen und somit die Unterschiede in den Off-Label-Tagen zwischen den Einrichtungen
und Diagnosen auf die unterschiedlichen Anwendungstage zurückführbar wären. Da
insbesondere im ambulanten Bereich die Anwendungsdauer im Mittel länger ist als
im stationären Bereich (siehe Anhang A.1 Abbildung
\ref{fig:Anwendungsdauer_ecdf}) würden wir nur anhand der durchschnittlichen
Off-Label-Tage, keine wertvolle Aussage über Unterschiede zwischen ambulanten und
stationären Einrichtungen im Off-Label-Use treffen können.\bigskip

\begin{figure}[ht]
    \centering
    \includegraphics[width = \textwidth]{../Images/OLU_Proportion_ecdf.png}
    \caption{Empirische Verteilungsfunktion der OLU-Proportion auf Patientenebene.}
    \label{fig:OLU_Proportion_ecdf.png}
\end{figure}

Wir wollen die OLU-Proportion der Patienten modellieren und untersuchen, ob die
Einrichtung und Diagnose einen Einfluss auf die OLU-Proportion haben. In
Abbildung \ref{fig:OLU_Proportion_ecdf.png} ist die Verteilung der
OLU-Proportionen dargestellt. Wir beobachten in der Abbildung, dass die
Verteilung der OLU-Proportion der Patienten in der Stichprobe
zwei große Sprungstellen aufweist. Die erste Sprungstelle befindet sich bei einer
OLU-Proportion von 0 und die zweite Sprungstelle befindet sich bei einer OLU-Proportion von 1. Diese
Sprungstellen resultieren daraus, dass einige Patienten entweder keine
Off-Label-Verordnungen erhalten haben oder ausschließlich Off-Label-Verordnungen erhalten
haben. Insbesondere entstehen die Sprungstellen nicht ausschließlich dadurch,
dass Patienten nur eine Verordnung erhalten haben, was dadurch ersichtlich wird,
dass Patienten mit einer OLU-Proportion von 0 im Mittel 5 Verordnungen erhalten haben
und Patienten mit einer OLU-Proportion von 1 im Mittel 3 Verordnungen erhalten haben
(siehe Anhang B OLU\_Proportion.Rmd Mittlere Anzahl der Verordnungen). Wir
schließen daraus, dass die Sprungstellen keine Artefakte der Datenerhebung sind,
sondern tatsächlich aufgrund der Verordnungspraxis entstehen und somit bei der
Analyse berücksichtigt werden müssen.

\subsection{Modellklasse}
\label{oluprop_model_class}

Wir beobachten weiter, dass die Verteilung der OLU-Proportion der Patienten im
Intervall $(0, 1)$ unimodal und rechtsschief ist und viele Ausprägungen
aufweist. Dies suggeriert, dass wir die OLU-Proportion der Patienten im
Intervall $(0, 1)$ der Patienten als stetige Variable auffassen können, die eine
unimodale, rechtsschiefe Verteilung aufweist. Eine mögliche Verteilungsklasse
für die OLU-Proportion der Patienten im Intervall $(0, 1)$ ist die
Beta-Verteilung, die auf dem Intervall $(0, 1)$ definiert ist und genügend
Flexibilität aufweist, um die beobachtete Verteilung der OLU-Proportion der
Patienten auf dem Intervall $(0, 1)$ zu modellieren. Die Beta-Verteilung ist
jedoch nicht im Stande die Sprungstellen bei einer OLU-Proportion von $0$ und
$1$ zu modellieren. Daher unterteilen wir die OLU-Proportion der Patienten in
drei Klassen: Patienten mit einer OLU-Proportion von $0$,
Patienten mit einer OLU-Proportion im Intervall $(0, 1)$ und Patienten mit einer
OLU-Proportion von $1$. Um diese drei Klassen zu modellieren, schlagen wir vor,
die OLU-Proportion der Patienten mit einer Beta-0-1-Inflated-Verteilung zu modellieren, siehe
\citet{rigbyDistributionsModellingLocation}. Die
Beta-0-1-Inflated-Verteilung ist eine Mischverteilung aus drei Komponenten:
einem diskreten Wert $0$ mit Wahrscheinlichkeit $p_0$, einem diskreten Wert $1$
mit Wahrscheinlichkeit $p_1$ und einer Beta-Verteilung $BE(\mu,\sigma)$ im
Intervall $(0, 1)$ mit Wahrcheinlichkeit $(1 - p_0 - p_1)$. Die
Beta-0-1-Inflated-Verteilung ist in der Lage, die beobachtete Verteilung der
OLU-Proportion der Patienten zu modellieren und die Sprungstellen bei einer
OLU-Proportion von $0$ und $1$ zu berücksichtigen. Im Folgenden werden wir die
Beta-0-1-Inflated-Verteilung genauer beschreiben.\bigskip

Die Wahrscheinlichkeitsdichte der Beta-0-1-Inflated-Verteilung 
$BEINF(\mu, \sigma, \nu, \tau)$, ist definiert als
\begin{equation}
    \label{eq:beinf}
    f_Y(y\mid \mu, \sigma, \nu, \tau) = \begin{cases}
        p_0 & \text{für } y = 0, \\
        (1 - p_0 - p_1) \cdot f_{W}(y \mid \mu, \sigma) & \text{für } y \in (0, 1),\\
        p_1 & \text{für } y = 1
    \end{cases}
\end{equation}
wobei $f_{W}(y \mid \mu, \sigma)$ die Dichte der Beta-Verteilung ist, also ${W \sim BE(\mu, \sigma)}$, mit 
$0 < \mu < 1$ und $0 < \sigma < 1$ und $p_0 = \frac{\nu}{1 + \nu + \tau}$ und $p_1 = \frac{\tau}{1 + \nu + \tau}$.

Durch Umformung obiger Gleichung für $p_0$ und $p_1$ erhalten wir
\begin{equation}
    \label{eq:beinf_nutau}
    \nu = \frac{p_0}{1 - p_0 - p_1} \quad \text{und} \quad
    \tau = \frac{p_1}{1 - p_0 - p_1}.
\end{equation}

Die Parameter $\nu$ und $\tau$ können als Chancenverhältnisse interpretiert
werden. Das bedeutet, dass $\nu$ die Chancen von $y = 0$ gegenüber $0 < y < 1$
beschreibt und $\tau$ die Chancen von $y = 1$ gegenüber $0 < y < 1$ beschreibt.
Sei also beispielsweise $\nu = \tau = 0.5$ und $\P(0 < Y < 1) = 0.5$, dann ist
$p_0 = p_1 = 0.5 \cdot \P(0 < Y < 1) = 0.25$. Eine solche Reparametrisierung hat
den Vorteil, dass die Parameter $\nu$ und $\tau$ unabhängig voneinander
interpretierbar sind. Das ist bei $p_0, p_1 \text{ und } (1-p_0-p_1)$ nicht der
Fall, da diese sich zu $1$ addieren müssen.



\bigskip

Die Parameter $\mu$ und $\sigma$ der Beta-Verteilung $BE(\mu, \sigma)$
beschreiben die Lage und Streuung der Beta-Verteilung im Intervall $(0, 1)$ im
folgenden Sinne, siehe \citet{rigbyDistributionsModellingLocation}:
\begin{align*}
    \label{eq:be_mu_sigma}
    \E(W) &= \mu\\
    \V(W) &= \mu (1-\mu) \sigma^2
\end{align*}

Aus den obigen Definitionen und Gleichungen folgt für
den Erwartungswert
\begin{equation}
    \label{eq:beinf_exp}
    \E(Y) = \frac{\mu + \tau}{1 + \nu + \tau}.
\end{equation}

Wir sehen anhand von Gleichung \ref{eq:beinf_exp}, dass der Erwartungswert durch
die Parameter $\mu, \nu$ und $\tau$ beeinflusst wird, aber nicht durch den
Parameter $\sigma$. Der Erwartungswert der Beta-0-1-Inflated-Verteilung ist also
unabhängig von der Streuung der Beta-Verteilung im Intervall $(0, 1)$. Dies ist
ein Vorteil der Beta-0-1-Inflated-Verteilung, da wir den Erwartungswert der
OLU-Proportion der Patienten modellieren wollen, aber nicht die Streuung der
OLU-Proportion der Patienten.\bigskip

Wir können außerdem anhand von Gleichung \ref{eq:beinf_exp} sehen, dass der
Erwartungswert der Beta-0-1-Inflated-Verteilung mit größer werdenden $\mu$
und $\tau$ steigt und mit größer werdenden $\nu$ sinkt.\bigskip

\subsection{Modelle}

Im Folgenden modellieren wir die OLU-Proportion der Patienten bedingt auf deren
Einrichtung, Diagnose, Geschlecht und Alter, mittels einer
Beta-0-1-Inflated-Verteilung. Dazu schätzen wir die Parameter
$\mu, \sigma, \nu \text{ und } \tau$ der Beta-0-1-Inflated-Verteilung.\bigskip

Dazu definieren wir den linearen Prädiktor $\eta^{(\theta)}$ als
\begin{align}
    \label{eq:eta}
    \begin{split}
    \eta^{(\theta)} := & \bx^\top \bbeta^{(\theta)}\\
    =& \beta_0^{(\theta)} + \beta_{1,1}^{(\theta)} \cdot \text{Einrichtung}_1 +
    \dots +  \beta_{1,K-1}^{(\theta)} \cdot \text{Einrichtung}_{K-1} +\\
    & \beta_{2}^{(\theta)} \cdot \text{Diagnose}_{\text{maligne}} + \beta_{3}^{(\theta)} \cdot \text{Diagnose}_{\text{neurologisch}} +\\
    & \beta_4^{(\theta)} \cdot \text{Geschlecht}_w + 
    \beta_5^{(\theta)} \cdot \text{Alter},
    \end{split}
\end{align}

wobei $\theta \in \{\mu , \nu, \tau\}$. Den Parameter $\sigma$ nehmen wir als
konstant an, d.h. der lineare Prädiktor für $\sigma$ ist $\eta^{(\sigma)} =
\beta_0^{(\sigma)}$. Das Merkmal Einrichtung ist eine der vier
Einrichtungsunterteilungen, d.h. Einrichtung2, Einrichtung3, Einrichtung13 oder
Einrichtung16 und $\text{Einrichtung}_{i}$ ist eine Indikatorvariable, die für
die möglichen Ausprägungen von den Einrichtungsunterteilungen steht. Wir verwenden als
Referenzkategorie bei Einrichtung2 ambulante Einrichtungen, bei Einrichtung3
SAPV und bei Einrichtung13 und Einrichtung16 die Einrichtung
SAPV\_D. Für das Merkmal Diagnose verwenden wir als Referenzkategorie die
Diagnose internistisch und für das Merkmal Geschlecht das Geschlecht
männlich. Das Merkmal Alter ist zentriert, d.h. wir haben den Mittelwert des
Alters von allen Patienten abgezogen.


Die Standard-Linkfunktionen, die die Parameter $\mu, \nu, \tau \text{ und }
\sigma$ zu den linearen Prädiktoren in Beziehung setzen
(\citet{rigbyDistributionsModellingLocation}), sind
\begin{align}
    \label{eq:link}
    \begin{split}
    \eta^{(\mu)} &= \text{logit}(\mu) = \log(\frac{\mu}{1 - \mu}),\\
    \eta^{(\nu)} &= \log(\nu),\\
    \eta^{(\tau)} &= \log(\tau),\\
    \eta^{(\sigma)} &= \text{logit}(\sigma) = \log(\frac{\sigma}{1 - \sigma}).
    \end{split}
\end{align}

Die Schätzung der Parameter $\mu, \nu, \tau \text{ und } \sigma$ erfolgt in R
mittels der Funktion \texttt{gamlss} aus dem Paket \texttt{gamlss}, siehe
\citet{stasinopoulosGeneralizedAdditiveModels2007}. Die Schätzung erfolgt
mittels penalisierter Maximum-Likelihood-Schätzung, siehe
\citet{stasinopoulosFlexibleRegressionSmoothing2017}: S. 62.\bigskip

Wir generieren insgesamt vier Modelle, die die OLU-Proportion der Patienten
modellieren. Die Verteilungsannahme ist in allen Modellen für alle Patienten $p,
p = 1,\dots,P$

\begin{equation}
    \label{eq:oluprop_dist}
    \text{OLU-prop}_p\mid \bx_p \sim BEINF(\mu_p, \sigma, \nu_p, \tau_p).
\end{equation}

Die resultierenden Modelle bezeichnen wir als \textbf{Modell2, Modell3,
Modell13} und \textbf{Modell16}, bzw. im Allgemeinen als \textbf{ModellK}
basierend auf der Unterteilung der Einrichtungen, die in den Modellen
berücksichtigt wird, siehe Abschnitt \ref{data_patient}. Insgesamt schätzen wir
bei ModellK $3 \cdot (K + 4) + 1$ Koeffizienten. Für Modell2 sind das insgesamt
$3 \cdot (2 + 4) + 1 = 21$ Koeffizienten und für Modell16 sind das insgesamt $3
\cdot (16 + 4) + 1 = 61$ Koeffizienten.

\subsection{Ergebnisse}
\label{oluprop_results}

Mit den oben definierten Modellen untersuchen wir in Abschnitt
\ref{oluprop_var_assoc}, ob die OLU-Proportion der Patienten mit
der Einrichtung oder Diagnose assoziiert ist. In Abschnitt
\ref{oluprop_model_fit} untersuchen wir die Modellgüte, um unsere Annahmen über
die Verteilung zu überprüfen. Anschließend
geben wir in Abschnitt \ref{oluprop_model_sum} eine Zusammenfassung der Ergebnisse des Modells für die Parameter
$\mu, \nu, \tau \text{ und } \sigma$. In Abschnitt
\ref{oluprop_erwartete_off_label_prop} berechnen wir die bedingte erwartete Verteilung der
OLU-Proportion der Patienten für die verschiedenen
Einrichtungen und Diagnosen und vergleichen diese miteinander.

\subsubsection{Assoziation der Variablen}
\label{oluprop_var_assoc}

Um zu untersuchen, ob diese mit der OLU-Proportion der Patienten assoziiert
sind, definieren wir die Nullhypothesen 
\begin{align}
    \label{eq:h0}
    H_0^{E} & : \forall \theta \in \{\mu, \nu, \tau\} \text{ gilt } \beta_{1,1}^{(\theta)} = \dots = \beta_{1,K-1}^{(\theta)} = 0,\\
    H_0^{D} & : \forall  \theta \in \{\mu, \nu, \tau\} \text{ gilt } \beta_{2}^{(\theta)} = \beta_{3}^{(\theta)} = 0.
\end{align}

Wir testen die Nullhypothesen $H_0^{E}$ und $H_0^{D}$ für jedes der vier Modelle
mittels Likelihood-Quotienten-Tests. Wir nutzen hierzu die in R für
\texttt{gamlss}-Modelle vorhandene Implementierung \texttt{LR.test}, siehe
\citet{stasinopoulosFlexibleRegressionSmoothing2017}. Hierzu vergleichen wir
ModellK mit einem reduzierten Modell, bei dem die Koeffizienten für die
Einrichtung (unter $H_0^{E}$) bzw. Diagnose (unter $H_0^{D}$) auf $0$ gesetzt
sind. Wir bezeichnen das reduzierte Modell unter $H_0^{E}$ als
\textbf{ModellK-E} und das reduzierte Modell unter $H_0^{D}$ als
\textbf{ModellK-D}. Die Ergebnisse der Likelihood-Quotienten-Tests sind in
Tabelle \ref{table:oluprop_var_assoc} dargestellt.

\begin{table}[ht]
    \begin{center}
    \begin{tabular}{|l|l|r|r|r|}
        \hline
        \textbf{Modell} & \textbf{reduziertes Modell} & \textbf{df} & $\mathbf{\chi^2}$ & \textbf{$p$-Wert}\\
        \hline
        \multirow{2}{4em}{Modell2} & Modell2-E & 3	& 1.99 & 0.575\\
        & Modell2-D & 6	& 12.94 & * 0.044 \\
        \hline
        \multirow{2}{4em}{Modell3} & Modell3-E & 6	& 3.55 & 0.74\\
        & Modell3-D & 6	& 12.84 & * 0.046 \\
        \hline
        \multirow{2}{4em}{Modell13} & Modell13-E & 36	& 77.24 & *** 0.000\\
        & Modell13-D & 6	& 12.29 & 0.056\\
        \hline
        \multirow{2}{4em}{Modell16} & Modell16-E & 45	& 85.25 & *** 0.000\\
        & Modell16-D & 6	& 12.09 & 0.060\\
        \hline
        \multicolumn{5}{| l |}{Signifikanzniveaus: '*' $p < 0.05$, '**' $p < 0.01$, '***' $p < 0.001$}\\ [1ex] 
        \hline
    \end{tabular}
    \caption{Ergebnisse des Likelihood-Quotienten-Tests für die Assoziation der Einrichtung und Diagnose mit der OLU-Proportion der Patienten.}
    \label{table:oluprop_var_assoc}
    \end{center}
\end{table}

Wir sehen anhand der Ergebnisse in Tabelle \ref{table:oluprop_var_assoc}, dass
für Modell2 und Modell3 die Diagnose statistisch signifikant auf dem 5\% Niveau
mit der OLU-Proportion der Patienten assoziiert ist, d.h. die Nullhypothese
$H_0^{D}$ wird abgelehnt. Wir sehen auch,
dass der $p$-Wert für ModellK-D für alle Tests kleiner als $0.1$ ist. Das
legt nahe, dass die Diagnose einen Einfluss auf die OLU-Proportion der Patienten
hat, dieser Einfluss jedoch nicht stark genug ist, um die Nullhypothese bei allen Modellen auf dem 5\% Niveau abzulehnen.
Für Modell13 und Modell16 ist die Einrichtung
statistisch signifikant auf dem 5\% Niveau mit der OLU-Proportion der Patienten
assoziiert, d.h. die Nullhypothese $H_0^{E}$ wird abgelehnt. Im Gegenzug dazu spricht der hohe $p$-Wert für ModellK-E bei Modell2
und Modell3 dafür, dass $H_0^{E}$ nicht abgelehnt wird und das Setting gemäß Einrichtung2 und Einrichtung3
keinen Einfluss auf die OLU-Proportion der Patienten hat, jedoch die unterschiedlichen 
Einrichtungen gemäß Einrichtung13 und Einrichtung16 sich stark voneinander unterscheiden.

\subsubsection{Modellgüte}
\label{oluprop_model_fit}

Um die Modellgüte bzw. die Modellannahmen zu überprüfen, benutzen wir
Quantil-Residuen, siehe \citet{dunnRandomizedQuantileResiduals1996} oder
\citet{stasinopoulosFlexibleRegressionSmoothing2017} für eine mathematische
Beschreibung der Quantil-Residuen. Die Quantil-Residuen dienen dazu, die
empirischen Quantile der OLU-Proportion der Patienten $\text{OLU\_prop}_p$ mit
den Quantilen der geschätzten Verteilung der OLU-Proportion der Patienten
$\widehat{\text{OLU\_prop}_p}$, sowie mit den erwarteten OLU-Proportionen der
Patienten $\E(\text{OLU\_prop}_p)$ zu vergleichen.
% Wir verwenden zur Darstellung der Quantil-Residuen die Funktion \texttt{plot.gamlss} aus dem Paket \texttt{gamlss}, siehe \citet{stasinopoulosFlexibleRegressionSmoothing2017}.
Quantil-Residual-Plots sind ähnlich zu Residual-Plots für lineare Modelle zu
interpretieren. Wir erwarten, dass die Quantil-Residuen um die Nulllinie streuen
und keine systematischen Muster aufweisen. Außerdem erwarten wir, dass die
Quantil-Residuen normalverteilt sind, was wir anhand der Dichtefunktion der
Quantil-Residuen und des QQ-Plots der Quantil-Residuen gegen die Quantile der
Standardnormalverteilung überprüfen können.

\begin{figure}[H]
    \centering
    \includegraphics[width = \textwidth]{../Images/OLU_Proportion_Quantile_Residuals16.png}
    \caption{Quantil-Residuen-Plots für Modell16.}
    \label{fig:OLU_Proportion_Quantile_Residuals16}
\end{figure}

In Abbildung \ref{fig:OLU_Proportion_Quantile_Residuals16} sehen wir die
Quantil-Residual-Plots für Modell16. Im linken Panel sehen wir die
Quantil-Residuen gegen die geschätzten erwarteten OLU-Proportionen der
Patienten $\widehat{\text{OLU\_prop}_p}$. Hier sehen wir, dass die
Quantil-Residuen um die Nulllinie streuen und keine Unterschiede in der
Streuung der Quantil-Residuen für die verschiedenen OLU-Proportionen der
Patienten aufweisen. Das deutet darauf hin, dass unsere Quantil-Residuen
unabhängig von der geschätzten OLU-Proportion der Patienten sind. Im mittleren
Panel sehen wir die geschätzte Dichtefunktion der Quantil-Residuen (in schwarz)
und die Dichtefunktion der Standardnormalverteilung (in rot). Hier sehen wir, dass die Dichtefunktion der Quantil-Residuen der Standardnormalverteilung sehr nahe kommt, was darauf hindeutet, dass unsere
Quantil-Residuen normalverteilt sind. 
Im rechten Panel wird der QQ-Plot dargestellt, in dem die Quantil-Residuen gegen die theoretischen Quantile einer Standardnormalverteilung geplottet sind. Die Datenpunkte liegen eng entlang der 45-Grad-Referenzlinie (in rot), was ebenfalls darauf hindeutet, dass die Verteilung der Quantil-Residuen eine gute Übereinstimmung mit der Standardnormalverteilung zeigt. Dies bestätigt visuell die Annahme einer Normalverteilung der Residuen.
Die Ergebnisse der Quantil-Residual-Plots von Modell2,
Modell3 und Modell13 sind in den Abbildungen
\ref{fig:OLU_Proportion_Quantile_Residuals2},
\ref{fig:OLU_Proportion_Quantile_Residuals3} und
\ref{fig:OLU_Proportion_Quantile_Residuals13} im Anhang A.1 dargestellt.\bigskip

Zusammenfassend können wir sagen, dass unsere Modelle die OLU-Proportion der
Patienten gut modellieren und die Verteilung der OLU-Proportion der Patienten
angemessen abbilden. Das stützt unsere Annahmen in Abschnitt
\ref{oluprop_model_class}, dass die Beta-0-1-Inflated-Verteilung die
OLU-Proportion der Patienten gut modelliert.

\subsubsection{Modellsummary}
\label{oluprop_model_sum}

In diesem Abschnitt schauen wir uns die Koeffizienten, die wir für unsere vier
Modelle erhalten, genauer an. Dazu betrachten wir zunächst exemplarisch die
Ergebnisse von Modell16, siehe Tabelle \ref{table:modell16_summary}. Wir
erhalten für alle vier Parameter $\mu, \sigma, \nu \text{ und } \tau$ die
geschätzten Koeffizienten der jeweiligen Kovariablen, sowie in Klammern darunter
deren Standardabweichung. Mit der üblichen Notation, wie am Ende der Tabelle
vermerkt, sind neben den Koeffizienten auch deren Signifikanzniveaus vermerkt.
Da wir den Parameter $\sigma$ als konstant betrachten, ist dessen geschätzter
Koeffizient gesondert am Ende der Tabelle angegeben.\bigskip

In Modell16 beschreibt der Intercept den erwarteten Wert für $\eta^{(\theta)}$,
wenn der Patient in der Einrichtung \textit{SAPV\_D} ist, die Diagnose
\textit{internistisch} hat und der Patient \textit{männlich} und
durchschnittlichen Alters (ca. \textit{73 Jahre}) ist.

Mittels der Linkfunktionen aus den Gleichungen (\ref{eq:link}) können wir den
Wert von $\eta^{(\theta)}$ umrechnen und erhalten dadurch die entsprechenden
geschätzten erwarteten Werte der Parameter, die wir im folgenden als $\mu_0,
\sigma_0, \nu_0 \text{ und } \tau_0$ bezeichnen. Die geschätzte erwartete
OLU-Proportion eines Patienten mit den oben genannten Merkmalsausprägungen
bezeichnen wir als $\text{OLU\_prop}_0$ und diese berechnet sich als
$\text{OLU\_prop}_0 = \frac{\mu_0 + \tau_0}{1 + \nu_0 + \tau_0}$. Konkret ergibt
sich
\begin{align}
    \label{eq:intercept}
    \begin{split}
    \mu_0 &= \text{logit}^{-1}(\beta_0^{(\mu)}) = \frac{1}{1 + \exp(-\beta_0^{(\mu)})} = \frac{1}{1 + \exp(-(-0.43))} = 0.39 \\
    \nu_0 &= \log^{-1}(\beta_0^{(\nu)}) = \exp(\beta_0^{(\nu)}) = \exp(-1.86) = 0.16\\
    \tau_0 &= \log^{-1}(\beta_0^{(\tau)}) = \exp(\beta_0^{(\tau)}) = \exp(-2.52) = 0.08\\
    \text{OLU\_prop}_0 &= \frac{\mu_0 + \tau_0}{1 + \nu_0 + \tau_0} = 38.39\%.
    \end{split}
\end{align}

Das bedeutet, die geschätzte erwartete OLU-Proportion für einen Patienten mit
den oben genannten Merkmalsausprägungen beträgt gemäß unserem Modell $38.39\%$.
Die Wahrscheinlichkeit, dass dieser Patient eine OLU-Proportion von $0$ hat,
beträgt $\frac{\nu_0}{1 + \nu_0 + \tau_0} \approx 13\%$ und für eine OLU-Proportion von $1$
beträgt die Wahrscheinlichkeit $\frac{\tau_0}{1 + \nu_0 + \tau_0} \approx 7\%$.\bigskip

Die Interpretation der restlichen Koeffizienten der jeweiligen Parameter erfolgt
im Kontext ihrer Linkfunktion (\ref{eq:link}). Erklärungen zur Interpretation
der $\log$-Linkfunktion und der logit-Linkfunktion finden sich in \citet{fahrmeirRegression}.\bigskip

Wir sehen in Tabelle \ref{table:modell16_summary}, dass die Einrichtungen sich im Parameter $\mu$ nicht signifikant von
der Referenzeinrichtung \textit{SAPV\_D} unterscheiden. Die Diagnosegruppen
\textit{maligne} und \textit{neurologisch} unterscheiden sich im Parameter $\mu$
ebenfalls nicht signifikant von der Referenzdiagnose \textit{internistisch}. Das
Geschlecht und das Alter haben keinen signifikanten Einfluss auf den Parameter
$\mu$.\bigskip

Die Einrichtungen unterscheiden sich im Parameter $\nu$ ebenfalls nicht
signifikant von der Referenzeinrichtung \textit{SAPV\_D}. Die Diagnosegruppen
\textit{maligne} und \textit{neurologisch} unterscheiden sich im Parameter $\nu$
ebenfalls nicht signifikant von der Referenzdiagnose \textit{internistisch}. Das
Geschlecht hat ebenfalls keinen signifikanten Einfluss auf den Parameter $\nu$.
Das Alter hingegen hat einen stark signifikanten positiven Einfluss auf den
Parameter $\nu$. Das bedeutet, dass mit steigendem Alter um ein Jahr die Chancen
von $\text{OLU\_prop}_p = 0$ gegenüber $0 < \text{OLU\_prop}_p < 1$ um den Faktor
$\exp(0.04) \approx 1.04$, also ca. $4\%$ steigen, wenn alle anderen Variablen konstant gehalten werden.\bigskip

Die Einrichtungen unterscheiden sich im Parameter $\tau$ nicht signifikant von
der Referenzeinrichtung \textit{SAPV\_D}. Die Diagnosegruppe \textit{maligne}
hat ebenfalls keinen signifikanten Einfluss auf den Parameter $\tau$. Die
Diagnosegruppe \textit{neurologisch} hat hingegen einen leicht signifikanten
positiven Einfluss auf den Parameter $\tau$. Das bedeutet, dass bei Patienten
mit der Diagnosegruppe \textit{neurologisch} die Chancen für $\text{OLU\_prop}_p =
1$ gegenüber $0 < \text{OLU\_prop}_p < 1$ um den Faktor $\exp(1.25) \approx 3.5$
größer sind, als bei Patienten mit der Diagnosegruppe \textit{internistisch}
(ceteris paribus). Das Geschlecht und das Alter haben keinen signifikanten
Einfluss auf den Parameter $\tau$.\bigskip

In den Modellen Modell2, Modell3 und Modell13 sind die Ergebnisse ähnlich zu
denen von Modell16, siehe Tabelle \ref{table:modell2_summary},
\ref{table:modell3_summary} und \ref{table:modell13_summary} im Anhang A.1.
Insbesondere sind die Koeffizienten für die Diagnosen, das Geschlecht und das
Alter in Modell2, Modell3 und Modell13 sehr ähnlich zu denen aus
Modell16, was darauf hindeutet, dass die Schätzungen robust gegenüber der Wahl
der Einrichtungsunterteilung sind.

\begin{longtable}[width=0.5\textwidth]{lrrrrr}
\toprule
\textbf{Koeffizienten} & $\mu$ & $\nu$ & $\tau$ & \textbf{OLU\_prop} \\ 
(Linkfunktionen) & (logit) & (log) & (log) & (identity) \\
\midrule\addlinespace[2.5pt]
Intercept & * -0.429    & ** -1.86    & ** -2.524    & *** 0.384 \\ 
 & (0.198)    & (0.703)    & (0.841)    & (0.064) \\ 
Diagnose maligne & 0.089    & -0.209    & -0.447    & 0.014 \\ 
 & (0.104)    & (0.36)    & (0.404)    & (0.034) \\ 
Diagnose neurologisch & 0.043    & 0.031    & * 1.246    & 0.091 \\ 
 & (0.189)    & (0.64)    & (0.552)    & (0.08) \\ 
Alter & 0.000    & ** 0.035    & 0.024    & -0.001 \\ 
 & (0.003)    & (0.013)    & (0.013)    & (0.002) \\ 
Geschlecht weiblich & -0.019    & -0.249    & 0.327    & 0.022 \\ 
 & (0.076)    & (0.29)    & (0.331)    & (0.03) \\ 
Einrichtung SAPV\_A & 0.227    & 0.286    & 0.762    & 0.065 \\ 
 & (0.199)    & (0.707)    & (0.831)    & (0.068) \\ 
Einrichtung SAPV\_B & 0.269    & -1.135    & -0.711    & 0.075 \\ 
 & (0.243)    & (1.205)    & (1.285)    & (0.078) \\ 
Einrichtung SAPV\_CA & -0.325    & -0.149    & -0.673    & -0.08 \\ 
 & (0.22)    & (0.825)    & (1.065)    & (0.066) \\ 
Einrichtung SAPV\_CB & -0.315    & -0.222    & -1.094    & -0.082 \\ 
 & (0.224)    & (0.829)    & (1.28)    & (0.068) \\ 
Einrichtung SAPV\_CC & -0.438    & 0.492    & 0.137    & -0.102 \\ 
 & (0.251)    & (0.84)    & (1.084)    & (0.074) \\ 
Einrichtung SAPV\_CD & 0.062    & 1.286    & 0.091    & -0.08 \\ 
 & (0.317)    & (0.848)    & (1.325)    & (0.094) \\ 
Einrichtung SAPV\_E & -0.007    & -1.035    & -0.313    & 0.02 \\ 
 & (0.206)    & (0.869)    & (0.936)    & (0.068) \\ 
Einrichtung Klinik\_A & 0.244    & -14.352    & 1.384    & * 0.201 \\ 
 & (0.282)    & (1065.361)    & (0.979)    & (0.092) \\ 
Einrichtung Klinik\_B & -0.156    & 0.586    & 0.636    & -0.029 \\ 
 & (0.228)    & (0.762)    & (0.946)    & (0.083) \\ 
Einrichtung Klinik\_C & -0.079    & 0.532    & -0.818    & -0.067 \\ 
 & (0.227)    & (0.759)    & (1.278)    & (0.073) \\ 
Einrichtung Uni\_A & -0.385    & -0.201    & -12.77    & -0.112 \\ 
 & (0.24)    & (0.977)    & (476.445)    & (0.071) \\ 
Einrichtung Uni\_B & 0.245    & -1.062    & 1.5    & * 0.19 \\ 
 & (0.239)    & (1.202)    & (0.881)    & (0.088) \\ 
Einrichtung Uni\_C & 0.014    & 0.045    & -0.296    & -0.004 \\ 
 & (0.254)    & (0.986)    & (1.302)    & (0.089) \\ 
Einrichtung Uni\_D & 0.101    & -0.936    & 0.479    & 0.077 \\ 
 & (0.23)    & (1.208)    & (1.001)    & (0.083) \\ 
Einrichtung Uni\_E & -0.212    & 0.771    & -0.424    & -0.093 \\ 
 & (0.264)    & (0.855)    & (1.303)    & (0.084) \\ 
\hline
        \multicolumn{5}{ l }{Signifikanzniveaus: '*' $p < 0.05$, '**' $p < 0.01$, '***' $p < 0.001$}\\ [1ex] 
        \multicolumn{5}{ l }{$\sigma$ (logit): *** -0.365    (0.043)}\\ [1ex]
        \hline
\caption{
{Summary Modell16.}
}
 \\ 
\label{table:modell16_summary}
\end{longtable}
\newpage
\subsubsection{Erwartete Off-Label-Proportion}
\label{oluprop_erwartete_off_label_prop}

Mit den im Abschnitt \ref{oluprop_model_sum} geschätzten Koeffizienten für die
Parameter $\mu, \nu \text{ und } \tau$ können wir die erwartete OLU-Proportion
für verschiedene Einrichtungen, Diagnosen, Geschlechter und Alter schätzen. Wir
interessieren uns allerdings nicht nur für die erwartete OLU-Proportion, sondern
auch dafür ob Unterschiede statistisch signifikant sind.\bigskip

Dafür definieren wir im Folgenden die $\xi_0$ als die erwartete OLU-Proportion
der Patienten, die in der Referenzeinrichtung $K$ sind, die Diagnose
\textit{internistisch} haben, männlich sind und durchschnittlichen Alters sind.
Mit der Notation aus Gleichung \ref{eq:eta} können wir $\xi_0$ als 
$$\xi_0 := \E(\text{OLU\_prop}_p \mid \bx = (1, 0, \dots, 0))$$ schreiben. Für
andere Merkmalsausprägungen definieren wir die Differenzen der erwarteten
OLU-Proportionen der Patienten zu $\xi_0$ als
\begin{align}
    \begin{split}
    \xi_{k} :=& \E(\text{OLU\_prop}_p \mid \text{Einrichtung}_k = 1) - \xi_0, \quad k \in \{1, \dots, K-1\},\\
    \xi_{d} :=& \E(\text{OLU\_prop}_p \mid \text{Diagnose}_d = 1) - \xi_0, \qquad d \in \{\text{maligne}, \text{neurologisch}\},\\
    \xi_{w} :=& \E(\text{OLU\_prop}_p \mid \text{Geschlecht}_w = 1) - \xi_0,\\
    \xi_{a} :=& \E(\text{OLU\_prop}_p \mid \text{Alter} = 1) - \xi_0.
    \end{split}
    \label{eq:xi}
\end{align}
Wir bedingen immer nur auf eine geänderte Merkmalsausprägung im Vergleich zur
Referenz. Wir bezeichnen die Differenzen $\xi_k, \xi_d, \xi_w \text{ und }
\xi_a$, sowie $\xi_0$ im Allgemeinen als $\xi$ und spezifizieren dieses wenn
nötig. Im Allgemeinen interessieren wir uns für die Verteilung von $\hat{xi}$.
In Berechnung \ref{eq:intercept} haben wir bereits für Modell16 $\hat{\xi_0}$
geschätzt als 
$$\hat{\xi_0} = 38.39\%.$$ Mit den geschätzten Koeffizienten für $\mu, \nu
\text{ und } \tau$ können wir im Allgemeinen $\hat{\xi}$ berechnen, wie in
Abschnitt \ref{oluprop_model_sum} erläutert. Um jedoch auch die Unsicherheit der
geschätzten $\xi$ zu bestimmen, führen wir ein nicht-parametrisches
Bootstrap-Verfahren durch, siehe \citet{kauermannStatisticalFoundationsReasoning2021a}, 204ff.\bigskip

Im Folgenden beschreiben wir das Bootstrap-Verfahren. Hierzu führen wir die
Folgende Notation ein:
\begin{align}
    \begin{split}
    \bX &:= \text{Designmatrix der Kovariablen},\\
    \bx_p &:= \text{Vektor der Kovariablen des Patienten } p,\\
    \bY &:= \text{Vektor der OLU-Proportionen der Patienten},\\
    y_p &:= \text{OLU-Prop}_p.\\
    \end{split}
    \label{eq:xi_hat}
\end{align}

Das Bootstrap-Verfahren zur Schätzung der Verteilung von $\xi$ verläuft wie folgt:
\begin{enumerate}
    \item Wir ziehen mit Zurücklegen $B$ Stichproben der Größe $P$ aus unseren
    Daten $(\bX, \bY)$. Wir bezeichnen die gezogenen Stichproben als
    $(\bX^*, \bY^*)^1, \dots, (\bX^*, \bY^*)^B$ und die jeweiligen
    Beobachtungen des $b$-ten Bootstrap-Datensatzes als $(\bx_1^*, y_1^*)^b, \dots,
    (\bx_P^*, y_P^*)^b$.

    \item Für jedes der $B$ Bootstrap-Datensätze fitten wir das Modell und
    berechnen $\hat{\xi^*}^b, b = 1,\dots, B$ wie oben beschrieben. 

    \item Als Nächstes schätzen wir die Varianz $Var(\hat{\xi})$ als
        \begin{equation*}
            \widehat{Var(\hat{\xi)}} = \frac{1}{B-1} \sum_{b=1}^B (\hat{\xi^*}^b - \bar{\hat{\xi^*}})^2
            \quad \text{ und } \widehat{sd(\hat{\xi)}} = \sqrt{\widehat{Var(\hat{\xi)})}},
        \end{equation*}
    wobei $\bar{\hat{\xi^*}} = \frac{1}{B} \sum_{b=1}^B \hat{\xi^*}^b$.
    
\end{enumerate}

Es gilt im Allgemeinen, dass die empirische Verteilung von $\hat{\xi^*}$,
bezeichnet als $\hat{F}_{\xi, B}(\cdot)$ für eine große Anzahl von
Bootstrap-Stichproben $B$ gegen die Verteilung von $\hat{\xi}$, bezeichnet als
$F_{\xi}(\cdot)$, konvergiert. In Abbildung
\ref{fig:mittelwert_diff_diagnosen_neurologisch} ist als Beispiel
die Bootstrap-Verteilung für $\widehat{\xi_{\text{neurologisch}}^*}$ für
Modell16 dargestellt. Wir sehen, dass $\hat{F}_{\xi, B}(\cdot)$ approximativ
normalverteilt ist. Mit der Bootstrap-Stichprobe können wir die Standartabweichung
schätzen, die Konfidenzintervalle für $\xi$ berechnen und Hypothesentests
durchführen. Die Null-Hypothesen definieren wir als
\begin{align}
    H_0^{\xi}: \xi = 0,
    \label{eq:null_hyp}
\end{align}
welche wir mittels Gauss-Test auf dem 5\% Niveau testen. In Tabelle
\ref{table:modell16_summary} in der Spalte \text{OLU\_prop} sind die geschätzten
$\xi$ für Modell16 aufgeführt, sowie deren geschätzte Standardabweichung,
$\widehat{sd(\hat{\xi)}}$ in Klammern und deren Signifikanzniveaus mit Sternchen
vermerkt. Wir sehen beispielsweise, dass für $\hat{\xi_0} = 38.39\%$ wir eine
geschätzte Standardabweichung von $\widehat{sd(\hat{\xi_0)}} = 0.064$ haben und
$\hat{\xi_0}$ signifikant von $0$ verschieden ist. Ebenfalls können wir sehen,
dass $\widehat{\xi_{\text{neurologisch}}} = 0.091$ ist und
$\widehat{sd(\widehat{\xi_{\text{neurologisch}})}} = 0.08$ ist. Die
Nullhypothese $H_0^{\xi_{\text{neurologisch}}}$ wird nicht auf dem 5\% Niveau
abgelehnt, da $\widehat{\xi_{\text{neurologisch}}}$ nicht statistisch
signifikant von $0$ verschieden ist.\bigskip

Die geschätzten $\xi$, sowie deren Standardabweichung und
Signifikanzniveaus, für Modell2, Modell3 und Modell13
sind in den Tabellen \ref{table:modell2_summary},
\ref{table:modell3_summary} und \ref{table:modell13_summary} im Anhang A.1
aufgeführt.\bigskip

Um die Unterschiede in den OLU-Proportionen der Patienten zwischen den Diagnosen
und Einrichtungen zu visuell zu vergleichen, stellen wir in den Abbildungen
\ref{fig:effektplot16_Einrichtung} und \ref{fig:effektplot16_Diagnose} die
geschätzten OLU-Proportionen in den jeweiligen Einrichtungen bzw. Diagnosen,
sowie deren 95\% Konfidenzintervall unter Modell16 dar. Für die
Referenzkategorien wäre dies $\hat{\xi_0}$, für die Diagnosen \textit{maligne}
und \textit{internistisch} wären dies $\hat{\xi_d} + \hat{\xi_0}$ und für die
Einrichtungen $k, k = 1,\dots,K-1$ wären dies $\hat{\xi_k} + \hat{\xi_0}$. Wir
sehen, dass wir für die Diagnose \textit{neurologisch} ein höheres $\xi$ im
Vergleich zu den anderen Diagnosen haben, da wir jedoch deutlich weniger
Patienten mit der Diagnose \textit{neurologisch} beobachtet haben, ist die
Unsicherheit in der Schätzung von $\xi_{\text{neurologisch}}$ größer, was sich
in einem breiteren Konfidenzintervall widerspiegelt. Für die Einrichtungen
sehen wir, dass die OLU-Proportionen der Patienten stark von der Einrichtung
abhängen. Insbesondere sehen wir, dass die Einrichtungen \textit{Klinik\_A} und
\textit{Uni\_B} eine signifikant höhere OLU-Proportion der Patienten haben, als die
Referenzeinrichtung \textit{SAPV\_D}. Wir können außerdem erkennen, dass die
Einrichtungen aus der \textit{SAPV\_C}-Gruppe sehr homogen sind.

Die Effektplots für Modell2, Modell3 und Modell13 sind in den Abbildungen
\ref{fig:effektplot2_Einrichtung}, \ref{fig:effektplot2_Diagnose},
\ref{fig:effektplot3_Einrichtung}, \ref{fig:effektplot3_Diagnose},
\ref{fig:effektplot13_Einrichtung} und \ref{fig:effektplot13_Diagnose} im Anhang
A.1 dargestellt.\bigskip

Zusammenfassend können wir sehen, dass es zwischen den Diagnosen in allen
Modellen keine signifikanten Unterschiede in den erwarteten OLU-Proportionen der Patienten
gibt, wobei die Diagnose \textit{neurologisch} tendenziell eine höhere erwartete
OLU-Proportion der Patienten hat. Das Geschlecht und das Alter haben keinen
signifikanten Einfluss auf die erwartete OLU-Proportionen der Patienten.

Bei den Einrichtungen sehen wir, dass in Modell2 und Modell3 wir keine
signifikanten Unterschiede in den erwarteten OLU-Proportionen der Patienten
zwischen den Einrichtungen haben. In Modell13 und Modell16 sehen wir jedoch
teilweise große Unterschiede in den erwarteten OLU-Proportionen der Patienten
zwischen den Einrichtungen und dass die Einrichtungen \textit{Klinik\_A} und
\textit{Uni\_B} signifikant höhere OLU-Proportionen der Patienten haben, als die
Referenzeinrichtung \textit{SAPV\_D}. Diese Ergebnisse sind konsistent mit 
den Ergebnissen in Tabelle \ref{table:oluprop_var_assoc}.

\begin{figure}[H]
    \centering
    \includegraphics[width = 0.8\textwidth]{../Images/mittelwert_diff_diagnosen_neurologisch.png}
    \caption{Bootstrap-Verteilung für $\widehat{\xi_{\text{neurologisch}}}$ mittels 10000 Bootstrap-Stichproben. Die vertikalen
    Linien markieren das 2.5\% und 97.5\% Quantil der Bootstrap-Verteilung.}
    \label{fig:mittelwert_diff_diagnosen_neurologisch}
\end{figure}

\begin{figure}[H]
    \centering
    \includegraphics[width = 0.8\textwidth]{../Images/Effekt_Einrichtung_16.png}
    \caption{Geschätzte OLU-Proportion und deren 95\% Konfidenzintervall in den verschiedenen Einrichtungen gemäß Modell16.}
    \label{fig:effektplot16_Einrichtung}
\end{figure}

\begin{figure}[H]
    \centering
    \includegraphics[width = 0.8\textwidth]{../Images/Effekt_Diagnose_16.png}
    \caption{Geschätzte OLU-Proportion und deren 95\% Konfidenzintervall für verschiedene Diagnosen Einrichtungen gemäß Modell16.}
    \label{fig:effektplot16_Diagnose}
\end{figure}



\subsection{Kritische Diskussion des gewählten Modells}
\label{oluprop_model_crit}

Die Wahl der Beta-0-1-Inflated-Verteilung für die Modellierung der
OLU-Proportionen der Patienten wurde aufgrund ihrer Fähigkeit, sowohl die
stetigen OLU-Proportionen im Intervall $(0, 1)$ als auch die diskreten
Sprungstellen bei einer OLU-Proportion von $0$ und $1$ zu modellieren,
getroffen. Eine alternative dazu wäre, die OLU-Proportionen der Patienten deren
OLU-Proportionen $0$ oder $1$ sind, leicht zu korrigieren, um sie ins offene
Intervall $(0,1)$ zu verschieben und anschließend die Beta-Verteilung zu
verwenden. Eine solche Korrektur könnte durchgeführt werden indem wir bei
Patienten mit einer OLU-Proportion von $0$ den Zähler in Gleichung
\ref{eq:oluprop_patient} um $1$ erhöhen und bei Patienten mit einer OLU-Proportion von
$1$ den Zähler um $1$ verringern. Jedoch konnten wir zeigen, dass die daraus
resultierenden Daten sich nicht gut mit der Beta-Verteilung modellieren lassen
und sehr stark von den ursprünglichen Daten abweichen. Außerdem würde diese
Korrektur die natürlichen Sprungstellen in der Off-Label-Use-Praxis ignorieren.
Die Beta-0-1-Inflated-Verteilung bietet die Möglichkeit, die Sprungstellen
angemessen zu berücksichtigen und ist daher angemessener für die Modellierung
der OLU-Proportionen der Patienten.\bigskip

% Es ist wichtig anzumerken, dass die Sprungstellen in den OLU-Proportionen nicht auf Fehler in der Datenerhebung zurückzuführen sind, sondern auf tatsächliche Praktiken in der Medikamentenverordnung. Daher ist es notwendig, ein Modell zu wählen, das diese Sprungstellen angemessen berücksichtigt. Die Beta-0-1-Inflated-Verteilung bietet diese Möglichkeit und ist daher die geeignete Wahl für die Modellierung der OLU-Proportionen der Patienten in unserer Analyse.

Ein Problem mit der Wahl der Beta-0-1-Inflated-Verteilung ist, dass für einige
Einrichtungen, in denen es keine oder nur wenig Patienten mit einer
OLU-Proportion von $0$ oder $1$ gibt, die Schätzung von $\nu$ und $\tau$ nicht
robust ist. Zudem ist zu beachten, dass wir keine Interaktionen zwischen den
Einrichtungen und Diagnosen berücksichtigt haben, obwohl es möglich ist, dass
die OLU-Proportionen der Patienten durch die Interaktionen zwischen den
Einrichtungen und Diagnosen beeinflusst werden. Jedoch würden wir dafür in
Modell13 bzw. Modell16 insgesamt 108 bzw. 135 Koeffizienten schätzen müssen, was
zu einer sehr komplexen Modellierung führen würde. Für Modell2 und Modell3
haben wir die Interaktionen zwischen den Einrichtungen und Diagnosen
einerseits zwecks Vergleichbarkeit und andererseits aufgrund dessen, dass
deskripitive Analysen keine Hinweise auf Interaktionen gegeben haben, nicht
berücksichtigt.


